% =====================================================================
% chapters/04_experimental_setup.tex
% Chapter 4 — Experimental Setup
% Project: Vision-Based Running Technique Classification from In-the-Wild
%          Internet Videos using Markerless Pose Estimation and Deep Learning
%
% Required preamble packages (declare in main.tex):
%   \usepackage{booktabs}  \usepackage{array}  \usepackage{tabularx}
%   \usepackage{amsmath}
%
% This chapter specifies the experimental protocol only. No empirical results
% appear here; they are reported in Chapter 5. Training hyper-parameters are
% reported only when supported by experiment evidence rather than guessed.
% Architecture parameter counts are configuration facts and are stated where known.
% =====================================================================

\chapter{Experimental Setup}
\label{ch:experimental-setup}

This chapter specifies how the framework of Chapter~\ref{ch:methodology} is
evaluated: the fixed dataset split and evaluation protocol, the experimental
pipelines and their roles, the model training configuration, the evaluation
metrics, and the design of the ablation and extended analyses. No results are
reported here; all empirical outcomes appear in Chapter~\ref{ch:results}.

% ---------------------------------------------------------------------
\section{Dataset Split and Evaluation Protocol}
\label{sec:setup-split}

All experiments use a single \emph{fixed split by video identifier}. Each video
is assigned to exactly one partition --- train, validation, or test --- and all
running cycles extracted from that video inherit its partition. This guarantees
that no cycle from a test video is ever seen during training, preventing the
model from recognising a specific video rather than the technique and thereby
avoiding the optimistic bias that video-level leakage would introduce
. The split is generated once and reused unchanged by every
pipeline and every ablation, so that all configurations are compared on
identical videos.

The validation partition is used for model selection and for decision-threshold
tuning (Section~\ref{sec:setup-metrics} and
Section~\ref{sec:abl-threshold}); the test partition is held out and used only
for final reporting. The per-split counts at both the video and cycle level are
reported with the results in Section~\ref{sec:dataset-stats}, since they
describe the data on which the controlled benchmarks are computed.

% ---------------------------------------------------------------------
\section{Experimental Pipelines}
\label{sec:setup-pipelines}

Five pipelines are evaluated. They share every stage except cycle source and
representation, which is what allows their differences to be attributed to those
two factors. Table~\ref{tab:pipelines} summarises their design and role.

\begin{table}[htbp]
\centering
\caption{Experimental pipelines, their cycle source, representation, models, and intended role. Only the first two rows constitute the main controlled comparison; the image pipeline is an extended representation comparison and the last two are supplementary.}
\label{tab:pipelines}
\renewcommand{\arraystretch}{1.3}
\footnotesize % Giữ nguyên cỡ chữ nhỏ để bảng gọn gàng

% Cấu hình mới: cột 1 và 2 để 'l' (vì chữ ngắn), các cột còn lại dùng 'X' hoặc định dạng của 'X' để tự động xuống dòng.
\begin{tabularx}{\textwidth}{l l X X X}
\toprule
Pipeline & Cycle source & Representation & Model(s) & Role \\
\midrule
Handmade landmark & Manual segmentation & Landmark $+$ velocity & MLP, 1D-CNN, BiLSTM, CNN-BiLSTM & Upper bound (not deployable) \\
\addlinespace
Autocorr landmark & Autocorr.\ (\texttt{contact\_diff}) & Landmark $+$ velocity & MLP, 1D-CNN, BiLSTM, CNN-BiLSTM & Main automatic baseline \\
\addlinespace
Autocorr image & Autocorr.\ (\texttt{contact\_diff}) & RGB crop & MobileNetV2--LSTM & Representation comparison \\
\addlinespace
Boundary detector & Learned BiLSTM boundaries & Landmark $+$ velocity & BiLSTM, CNN-BiLSTM & Supplementary \\
\addlinespace
Raw sequence & None (whole video) & Full landmark sequence & GRU, BiLSTM & Supplementary / exploratory \\
\bottomrule
\end{tabularx}
\end{table}


% ---------------------------------------------------------------------
\section{Model Training Configuration}
\label{sec:setup-training}

All models are trained on the train partition, selected on the validation
partition, and evaluated once on the test partition. The pose-sequence models
operate on eight-phase landmark inputs of shape $(B,\,8,\,78)$, where
$F{=}78$ features comprise 13 selected joints $\times$ 3 coordinates $+$ 3
velocity components (Section~\ref{sec:method-phases}). The layer-by-layer
architectures of the four landmark classifiers are shown in
Figure~\ref{fig:model_architectures}.

The four landmark-sequence models share a common training schedule confirmed
from the experiment notebooks: \textbf{AdamW} optimiser with learning rate
$10^{-3}$ and weight decay $10^{-4}$; \textbf{OneCycleLR} scheduler with
\texttt{pct\_start}$=0.2$; \textbf{BCEWithLogitsLoss} with a class-ratio
\texttt{pos\_weight}; gradient clipping (max norm $1.0$); and early stopping
with patience $15$. Batch size is $16$, maximum epochs $80$, and random seed
$42$.

The image model uses eight RGB crops per detected cycle and a MobileNetV2--LSTM
architecture \cite{sandler2018mobilenetv2, donahue2015lrcn}.
Table~\ref{tab:model-config} summarises the per-model configuration.
The audited MobileNetV2--LSTM experiment uses learning rate $10^{-3}$, batch
size $8$, $40$ epochs, dropout $0.4$, two recurrent layers, weight decay
$10^{-4}$, and random seed $42$.


\begin{table}[htbp]
\centering
\caption{Model configuration confirmed from the experiment notebooks. Parameter
counts are taken from the exported model-ablation report. The four
landmark-sequence models share a common optimizer/scheduler/loss (see text);
per-model architecture details follow Figure~\ref{fig:model_architectures}.}
\label{tab:model-config}
\renewcommand{\arraystretch}{1.18}
\scriptsize
\begin{tabularx}{\textwidth}{
>{\raggedright\arraybackslash}p{0.16\textwidth}
>{\raggedright\arraybackslash}p{0.18\textwidth}
>{\raggedright\arraybackslash}p{0.22\textwidth}
>{\raggedright\arraybackslash}p{0.10\textwidth}
>{\raggedright\arraybackslash}X}
\toprule
Model & Pipeline use & Input & Params & Architecture notes \\
\midrule
MLP & Handmade \& autocorr landmark & $(B,8,78)$ & $193{,}025$ &
  Flatten $\to$ Linear(624$\to$256)+ReLU $\to$ Dropout(0.3) $\to$
  Linear(256$\to$128)+ReLU $\to$ Dropout(0.3) $\to$ Linear(128$\to$1). \\
\addlinespace
1D-CNN & Handmade \& autocorr landmark & $(B,8,78)$ & $40{,}257$ &
  Conv1d(78$\to$64, $k$=3)+ReLU+BN $\to$ Conv1d(64$\to$128, $k$=3)+ReLU+BN
  $\to$ AdaptiveAvgPool1d(1) $\to$ Dropout(0.3) $\to$ Linear(128$\to$1). \\
\addlinespace
BiLSTM & Handmade \& autocorr landmark & $(B,8,78)$ & $173{,}441$ &
  2-layer bidirectional LSTM ($h$=64) $\to$ concat[$h_\text{fwd}$, $h_\text{bwd}$]
  (128-d) $\to$ LayerNorm(128) $\to$ Dropout(0.3) $\to$ Linear(128$\to$1). \\
\addlinespace
CNN-BiLSTM & Handmade \& autocorr landmark & $(B,8,78)$ & $82{,}113$ &
  Conv1d(78$\to$64, $k$=3)+ReLU+BN $\to$ 1-layer BiLSTM ($h$=64) $\to$
  concat (128-d) $\to$ LayerNorm(128) $\to$ Dropout(0.3) $\to$ Linear(128$\to$1). \\
\addlinespace
MobileNetV2--LSTM & Autocorr image pipeline & 8$\times$RGB $224{\times}224$ & --- &
  MobileNetV2 backbone (frozen or fine-tuned) $\to$ 2-layer LSTM
  ($h$=256, dropout=0.4) $\to$ Linear($\to$1).
  lr=$10^{-3}$; batch=8; epochs=40; seed=42. \\
\addlinespace
GRU / BiLSTM (raw) & Supplementary raw-sequence & Full video landmark seq. & --- &
  Whole-video recurrent model; one sample per video. Treated as supplementary
  evidence; detailed hyper-parameters not audited. \\
\bottomrule
\end{tabularx}
\end{table}

Parameter counts for the four landmark models are taken directly from the
exported model-ablation report and are used as a reference for model complexity.
The MobileNetV2--LSTM and raw-sequence rows are included for completeness;
their parameter counts were not exported. All landmark classifiers were trained
under the shared schedule described above; the MobileNetV2--LSTM and
raw-sequence models follow separate configurations noted in their respective
rows.

% ---------------------------------------------------------------------
\section{Evaluation Metrics}
\label{sec:setup-metrics}

The following metrics are used.
\emph{Accuracy}, \emph{precision}, \emph{recall}, and \emph{F1} follow their standard definitions.
The \emph{macro-averaged F1} (macro-F1) is the unweighted mean of per-class F1 over the two classes (correct, false):
\begin{equation}
\mathrm{macro\text{-}F1} =
\frac{1}{2}\left(\mathrm{F1}_{\text{correct}} + \mathrm{F1}_{\text{false}}\right).
\end{equation}

\emph{ROC-AUC} is the area under the receiver-operating-characteristic curve,
which plots the true-positive rate against the false-positive rate as the
decision threshold varies; it summarises ranking quality independently of any
single threshold, with $0.5$ corresponding to chance \cite{fawcett2006roc}.

\paragraph{Why macro-F1 is the primary metric.}
Because the dataset is class-imbalanced at the video level (more false than
correct videos, Section~\ref{sec:dataset-stats}), accuracy can be misleading: a
classifier that favours the majority class can attain deceptively high accuracy
while failing on the minority class. Macro-F1 weights both classes equally
regardless of their frequency, so a model must perform well on \emph{both}
correct and false technique to score well \cite{sokolova2009metrics}. Macro-F1 is
therefore reported as the primary metric throughout, with accuracy, per-class
F1, and ROC-AUC reported alongside it for completeness. ROC-AUC is particularly
useful for separating ranking quality from threshold choice, which is examined
explicitly in Section~\ref{sec:abl-threshold}.

% ---------------------------------------------------------------------
\section{Ablation Study Design}
\label{sec:setup-ablations}

The ablations and extended analyses each vary one factor while holding the
remainder of the pipeline fixed, so that the measured effect can be attributed
to that factor. Table~\ref{tab:ablations} summarises the design; the
corresponding results appear in Chapter~\ref{ch:results}.

\begin{table}[htbp]
\centering
\caption{Design of the ablation studies and extended analyses. Each study varies
one factor on a fixed pipeline and split.}
\label{tab:ablations}
\renewcommand{\arraystretch}{1.3}
\footnotesize
\begin{tabularx}{\textwidth}{l X l X}
\toprule
Study & Varied factor & Pipeline & Held fixed \\
\midrule
Model ablation & Classifier (MLP / 1D-CNN / BiLSTM / CNN-BiLSTM) & Handmade \& Autocorr & Data, split, representation, threshold \\
Feature ablation & Feature set (selected, full-33, $+$velocity, $+$AP/V, $+$contact, $+$all) & Handmade & Model, split, frames \\
Signal ablation & Detection signal (head-Y, hip-Y, heel-Y, toe-Y, \texttt{contact\_diff}, \texttt{AP\_diff}, \texttt{V\_diff}) & Autocorr & Model, split, representation \\
Frames-per-cycle & Frames sampled per cycle (8 / 16 / 24 / 32) & Handmade & Model family, split, representation \\
Representation & Landmark vs.\ image & Autocorr (shared cycle source) & Cycle source, split \\
Cycle vs.\ video & Evaluation unit (cycle vs.\ video aggregation) & Autocorr (and others) & Model, split \\
Threshold & Decision threshold ($0.5$ vs.\ validation-tuned) & All & Model, split \\
Pose/cycle quality & Quality factor (visibility, autocorr., missing-ratio, duration, period) & Autocorr & Model, split \\
Fair-gap robustness & Cycle source under repeated folds and label permutations & Handmade vs.\ autocorr & MLP, video-level CV, threshold $0.5$ \\
\bottomrule
\end{tabularx}
\end{table}

\subsection{Model Ablation}
\label{sec:abl-model}
The four classifiers are compared on identical data for both the handmade and
autocorr pipelines, holding the representation, split, and decision threshold
fixed, to isolate the effect of architecture.

\subsection{Feature Ablation}
\label{sec:abl-feature}
On the handmade pipeline, the input feature set is varied --- from the compact
selected-landmark set through the full landmark set and several
derived-feature augmentations --- with the classifier and all other stages held
fixed, to measure the effect of feature representation.

\subsection{Signal Ablation}
\label{sec:abl-signal}
On the autocorr pipeline, the one-dimensional motion signal driving cycle
detection is varied across the seven candidates of
Table~\ref{tab:signals}, holding the classifier fixed. Both detection
statistics (coverage, period regularity, rejections) and downstream
classification are recorded, since signal choice affects both.

\subsection{Frames-per-Cycle Ablation}
\label{sec:abl-frames}
On the handmade pipeline, the number of frames sampled per cycle is varied
(8/16/24/32) for the sequence models, holding all else fixed, to test whether
temporal resolution beyond eight phases improves classification and at what
computational cost.

\subsection{Representation Ablation}
\label{sec:abl-representation}
The landmark and image representations are compared on the \emph{same}
autocorrelation cycle source and split, so that any difference is attributable
to representation rather than to cycle source or protocol. This analysis is kept
separate from the historical image-based development milestone, which was
recorded under a different protocol and is not a controlled benchmark.

\subsection{Cycle-Level versus Video-Level Evaluation}
\label{sec:abl-cycle-video}
Predictions are evaluated both per cycle and aggregated per video (majority
vote, mean probability, confidence-weighted vote), holding the model and split
fixed, to assess how aggregation affects measured performance. Video-level
results are treated as the more faithful estimate of practical performance.

\subsection{Threshold 0.5 versus Tuned Threshold}
\label{sec:abl-threshold}
For each model, performance at the default $0.5$ threshold is compared with
performance at a threshold selected on the \emph{validation} set (by maximising
validation macro-F1) and applied unchanged to test. Reporting both, together
with the change between them, exposes any threshold overfitting rather than
concealing it.

\subsection{Pose and Cycle Quality Analysis}
\label{sec:abl-quality}
Test cycles from the autocorr pipeline are grouped by pose- and cycle-quality
factors --- landmark visibility, autocorrelation strength, landmark
missing-ratio, cycle duration, and cycle period --- and classification
performance is compared across groups, to locate the dominant source of error in
the automatic pipeline. Because no canonical quality thresholds exist, groups are
formed by a median split, and the resulting group sizes are reported with the
results so that small groups are interpreted cautiously
(Section~\ref{sec:quality}).

\subsection{Fair Gap Robustness Check and Permutation Test}
\label{sec:abl-fair-gap}

The fixed split gives the primary controlled benchmark. Because the validation
and test partitions are small, an additional robustness check is used to test
whether the handmade--autocorr gap is specific to that split. Both cycle sources
are evaluated with the same MLP classifier, 5-fold video-level cross-validation,
the default threshold of $0.5$, and macro-F1 as the main metric. Splitting is
again performed by video identifier, so all cycles from one video remain within
the same fold.

A permutation test is also applied to the autocorr MLP setting. The observed
cross-validation macro-F1 is compared with a null distribution obtained by
randomly permuting the labels and rerunning the same cross-validation procedure.
The test is used as a diagnostic of separability under a simple classifier, not
as a replacement for the controlled benchmark. A non-significant result is
therefore interpreted conservatively: it indicates that the autocorr-derived
features are not robustly separable under that MLP configuration, rather than
proving that no useful signal exists.

% End of Chapter 4
